\documentclass[UTF8,twocolumn]{ctexart}
\usepackage[a4paper,margin=1.8cm]{geometry}
\usepackage{graphicx}
\usepackage{booktabs}
\usepackage{hyperref}
\usepackage{amsmath}

\title{温州话低资源大语言模型适配实验}
\author{Sidney Chai}
\date{\today}

\begin{document}
\maketitle

\begin{abstract}
本文基于多来源温州话语料，构建字符级预训练、指令微调和偏好优化三阶段实验，分析低资源方言场景下数据清洗、模型适配和评估方法的实际效果。
\end{abstract}

\section{研究目标}
说明温州话低资源 NLP 的背景、项目目标和成功标准。

\section{相关工作}
概述低资源语言处理、汉语方言计算、LoRA 微调和 DPO 偏好优化。

\section{数据集与预处理}
描述 13 个原始文件、勘误、去重、固定划分和防泄漏策略。

\section{实验一：字符级预训练}
报告模型配置、unigram baseline、loss 曲线和生成样例。

\section{实验二：SFT 翻译微调}
报告基座模型、LoRA 配置、训练设置和测试集结果。

\section{实验三：DPO 偏好优化}
报告偏好对构造、训练参数、SFT 与 DPO 对比，以及负结果可能原因。

\section{综合评估}
\input{figures/auto_metrics_table.tex}

\section{错误分析}
分析漏译、错译、幻觉、重复和输出语言错误。

\section{局限性}
讨论语料规模、方言书写差异、自动指标不足和 DPO 数据质量。

\section{结论}
总结三阶段实验的主要发现。

\bibliographystyle{plain}
\begin{thebibliography}{9}
\bibitem{vaswani2017attention} Vaswani et al. Attention Is All You Need. 2017.
\bibitem{hu2021lora} Hu et al. LoRA: Low-Rank Adaptation of Large Language Models. 2021.
\bibitem{rafailov2023dpo} Rafailov et al. Direct Preference Optimization. 2023.
\end{thebibliography}

\end{document}
